% ===========================================================================
% §1 — Introduction : l'objectif du projet
% ===========================================================================
\section{Introduction}
\label{sec:introduction}

Combien reste-t-il \emph{vraiment} à un ménage québécois, une fois payés ses
impôts et ses cotisations, et une fois reçus ses transferts --- allocations,
crédits et prestations ? La réponse est moins simple qu'il n'y paraît : au
Québec, le revenu disponible d'un ménage résulte de l'interaction d'une
vingtaine de programmes, administrés par deux ordres de gouvernement, chacun
avec ses propres seuils, taux, plafonds et règles de réduction.

Le ministère des Finances du Québec (MFQ) publie un calculateur officiel du
revenu disponible\footnote{\url{https://www.finances.gouv.qc.ca/ministere/outils_services/outils_calcul/revenu_disponible/outil_revenu.asp}}.
Cet outil donne le résultat, mais expose peu le \emph{raisonnement} : on y
voit les montants, pas les règles qui les produisent. Le projet
\emph{Revenu disponible --- Québec} part de l'intuition inverse : pour
comprendre la fiscalité, il faut pouvoir ouvrir la boîte noire.

\subsection{Objectif du projet}

Le projet poursuit quatre objectifs :

\begin{enumerate}
  \item \textbf{Reconstruire} le calcul du revenu disponible des ménages
        québécois, poste par poste, dans un code source lisible, commenté et
        publié (\ghurl), là où le calculateur officiel n'est disponible que
        sous forme compilée ;
  \item \textbf{Vérifier} chaque paramètre --- montant, taux, seuil ---
        contre la source officielle qui le fixe (Revenu Québec, Agence du
        revenu du Canada, Retraite Québec, Emploi et Développement social
        Canada) et citer la disposition législative correspondante ;
  \item \textbf{Valider} l'ensemble par comparaison systématique avec le
        calculateur officiel du MFQ : sur plusieurs milliers de cas types,
        les résultats concordent au cent près
        (section~\ref{sec:validation}) ;
  \item \textbf{Expliquer} : l'application Web rend chaque poste
        explorable --- description, objectif, règle de calcul, paramètres,
        références --- et permet de simuler des changements de paramètres,
        comme à l'annonce d'un budget.
\end{enumerate}

\subsection{À qui s'adresse ce document}

Ce guide vise d'abord les \textbf{étudiantes et étudiants en fiscalité}, qui
y trouveront, pour chaque poste, des règles de calcul rédigées en français
mais suffisamment précises pour être reproduites à la calculatrice --- elles
correspondent aux règles effectivement codées dans le moteur TypeScript,
y compris les cas limites (planchers, plafonds, arrondis) et les choix de
modélisation hérités du calculateur officiel. Le lecteur citoyen y trouvera
quant à lui une explication d'ensemble : pourquoi son taux marginal effectif
ne ressemble pas à son taux d'imposition affiché, d'où viennent les
\og trappes \fg{} à certains niveaux de revenu, et comment les programmes
s'emboîtent.

\subsection{Plan du document}

La section~\ref{sec:modele} délimite le modèle : ce qu'il couvre, ses
entrées, ses hypothèses. La section~\ref{sec:manuel} est le mode d'emploi de
l'application Web. La section~\ref{sec:architecture} décrit l'architecture du
calcul : l'ordre des postes et les bases de revenu intermédiaires qui les
relient. Les sections~\ref{sec:cotisations} à~\ref{sec:rd} détaillent les
vingt postes un à un. La section~\ref{sec:validation} explique la méthode de
validation et les limites du modèle.
